\documentclass[11pt]{article}
\usepackage[margin=1in]{geometry}
\usepackage{amsmath,amssymb,amsthm}
\usepackage{booktabs}
\usepackage{hyperref}

\newtheorem{theorem}{Theorem}
\newtheorem{proposition}{Proposition}
\newtheorem{lemma}{Lemma}
\newtheorem{conjecture}{Conjecture}
\newtheorem{remark}{Remark}
\newtheorem{definition}{Definition}
\newcommand{\mc}{\operatorname{MaxCut}}
\newcommand{\dd}{d}
\newcommand{\Opt}{\operatorname{Opt}}
\newcommand{\E}{\mathbb{E}}

\title{Optimal-Restriction Five-Sets and Global Cut Stability\\in Triangle-Free Graph Bipartization}
\author{Jinji Li}
\date{Version 3.0 draft -- September 2026}

\begin{document}
\maketitle

\begin{abstract}
For a graph $G$, let
\[
\dd(G)=|E(G)|-\mc(G)
\]
be the minimum number of edges whose deletion makes $G$ bipartite.  A
classical conjecture of Erd\H{o}s asserts that every triangle-free graph on
$N$ vertices satisfies $\dd(G)\le N^2/25$; for $N=5k$ this is
$\dd(G)\le k^2$.

We investigate a five-vertex induction target
\[
\exists Y\subseteq V(G),\ |Y|=5:\qquad
\dd(G)-\dd(G-Y)\le 2k-1.
\]
Earlier versions of this project emphasized local recoloring certificates.
The present note shifts to a global cut-stability phenomenon suggested by
exact computations.  For a maximum cut $f$ of $G$, let
$R_f(Y)$ be the number of monochromatic edges under $f$ incident with $Y$,
and let $\Delta_f(Y)$ be the additional decrease in monochromatic edges
obtained by reoptimizing the cut on $G-Y$.  Then exactly
\[
\dd(G)-\dd(G-Y)=R_f(Y)+\Delta_f(Y).
\]
We call $Y$ \emph{optimal-restriction preserving} if some maximum cut of
$G$ restricts to a maximum cut of $G-Y$, equivalently if
$\Delta_f(Y)=0$ for some $f\in\Opt(G)$.

On a saved canonical corpus of $19{,}270$ triangle-free graphs on
$15$ vertices, an exact full-cut audit finds that every saved graph admits
a five-set $Y$ with $\dd(G)-\dd(G-Y)\le5$ and with optimal-restriction
preservation.  The same conclusion holds for all $5{,}337$ saved graphs
with $L(G)\ge12$, all $2{,}103$ saved graphs with $\dd(G)>5$, and all
seven exactly identified edge-critical graphs.  No failure occurs in the
saved corpus.  A later optimal-face diagnostic finds at least $218$
low-defect preserving five-sets per saved graph.

We prove exact equivalences between optimal restriction, an intersection
of optimal-cut faces under deletion, and a five-vertex extension frontier.
Stronger shape properties of that frontier---monotonicity, one-step
Lipschitz behavior, and discrete convexity---are explicitly falsified by
saved witnesses.  Thus the surviving phenomenon is genuinely global: for
some five-set, the combined remainder-plus-extension cost is minimized at
an optimal remainder coloring even though the extension frontier itself
can oscillate strongly.

The full Erd\H{o}s conjecture remains open.  The contribution here is a
reproducible global-cut phenomenon, exact equivalences that isolate it,
and a sharply defined selection problem suggested by the full saved
$n=15$ corpus.
\end{abstract}

\section{Introduction}

For a finite simple graph $G$, define
\[
\dd(G)=\min\{|F|:F\subseteq E(G),\ G-F\text{ is bipartite}\}
      =|E(G)|-\mc(G).
\]
Erd\H{o}s conjectured that every triangle-free graph on $N$ vertices
satisfies
\[
\dd(G)\le \frac{N^2}{25}.
\]
The balanced blow-up of $C_5$ shows that $1/25$ would be sharp.  Balogh,
Clemen and Lidick\'y proved the global $N^2/23.5$ bound and the sharp
$N^2/25$ bound in two density ranges \cite{BCL}.  Ferudun gave a
computer-assisted proof of the sharp value for multiples of five through
$N=200$ \cite{Ferudun}.  The computations below are therefore used for
structure discovery and exact diagnostics, not as new finite extremal
evaluations.

The induction target motivating this project is the following.

\begin{conjecture}[Five-set induction target]
For every triangle-free graph $G$ on $5k$ vertices, there exists a five-set
$Y$ such that
\[
\dd(G)-\dd(G-Y)\le 2k-1.
\]
\end{conjecture}

If this holds for every $k$, then induction gives
\[
\dd(G)\le 1+3+\cdots +(2k-1)=k^2.
\]
The challenge is that deleting $Y$ may cause the remaining graph to choose
a substantially better global cut.  This note isolates that reoptimization
term and studies an empirical phenomenon in which it vanishes for a
suitable five-set.

\section{Corrected automatic threshold}

For a five-set $Y$, define
\[
\ell(Y)=2\dd(G[Y])+e(Y,G-Y),
\qquad
L(G)=\min_{|Y|=5}\ell(Y).
\]
The coordinated-flip selection inequality from the earlier project gives
\[
\dd(G)-\dd(G-Y)
\le
\left\lfloor
\dd(G[Y])+\frac{e(Y,G-Y)}2-\Phi_Y
\right\rfloor,
\]
where $\Phi_Y\ge0$ is an averaged recoloring gain.  Taking the empty flip
immediately yields
\[
\dd(G)-\dd(G-Y)
\le
\left\lfloor\frac{\ell(Y)}2\right\rfloor.
\]
Hence:

\begin{proposition}[Automatic threshold]
If $G$ has $5k$ vertices and some five-set $Y$ satisfies
\[
\ell(Y)\le 4k-1,
\]
then
\[
\dd(G)-\dd(G-Y)\le 2k-1.
\]
\end{proposition}

At $k=3$, the automatic range extends through $L(G)=11$.  This corrects
the threshold used in Version 1 of the project.  Consequently the older
$L=11$ exchange audit remains a valid structural computation but is not
needed to establish the induction inequality in that regime.

\section{Global reoptimization after deleting five vertices}

Fix a maximum cut $f$ of $G$, represented as a 2-coloring.  For a five-set
$Y$, put $H=G-Y$ and let $\operatorname{mono}_f(H)$ denote the number of
monochromatic edges of $H$ under the restriction of $f$.

Define
\[
R_f(Y)=\dd(G)-\operatorname{mono}_f(H),
\]
which is exactly the number of monochromatic edges of $f$ having at least
one endpoint in $Y$.  Define the remainder reoptimization gain
\[
\Delta_f(Y)=\operatorname{mono}_f(H)-\dd(H)\ge0.
\]
Then:

\begin{lemma}[Exact reoptimization decomposition]
For every maximum cut $f$ of $G$ and every five-set $Y$,
\[
\boxed{
\dd(G)-\dd(G-Y)=R_f(Y)+\Delta_f(Y).
}
\]
\end{lemma}

\begin{proof}
Since $f$ is globally optimal, $\dd(G)$ is its monochromatic-edge count.
Partition those monochromatic edges into the ones incident with $Y$ and
the ones wholly inside $H$.  The latter number is
$\operatorname{mono}_f(H)=\dd(H)+\Delta_f(Y)$.
\end{proof}

The local term $R_f(Y)$ is easy to average.  The difficulty is the global
term $\Delta_f(Y)$, which measures how much the entire remainder can
reorganize after deletion.

\section{Optimal-restriction five-sets}

\begin{definition}
A pair $(f,Y)$ is \emph{optimal-restriction preserving} if
$f\in\Opt(G)$ and the restriction $f|_{G-Y}$ is optimal for $G-Y$.
Equivalently,
\[
\Delta_f(Y)=0.
\]
We call $Y$ an \emph{H0 witness} if it is good for the induction bound and
is optimal-restriction preserving for at least one maximum cut:
\[
\dd(G)-\dd(G-Y)\le 2k-1,
\qquad
\exists f\in\Opt(G):\Delta_f(Y)=0.
\]
\end{definition}

For such a pair,
\[
\dd(G)-\dd(G-Y)=R_f(Y),
\]
so the global reoptimization difficulty disappears completely.

The following exact audit motivated the rest of the note.

\begin{proposition}[Saved-corpus H0 phenomenon]
In the saved canonical $n=15$ corpus used by this project, every one of
$19{,}270$ graphs has an H0 witness.  In particular, the property holds
for all $5{,}337$ saved graphs with $L(G)\ge12$, all $2{,}103$ saved graphs
with $\dd(G)>5$, and all seven exactly identified edge-critical graphs.
No failure occurs in the saved corpus.
\end{proposition}

This proposition is a statement about the precisely defined saved corpus.
It is not claimed that the saved corpus exhausts all triangle-free graphs
on $15$ vertices.

A representative-sample audit first suggested the phenomenon on $40$
graphs.  The exact full-corpus run then verified all $19{,}270$ saved
graphs, with independent recomputation of the optimum of every saved
witness core.  Both
\[
\max_G\min_{Y\text{ good}}\Delta_{\min}(Y)=0
\]
and the corresponding minimum cut-distance statistic were zero.

\section{Abundance rather than uniqueness}

The H0 witness is not rare in the saved data.  Subsequent diagnostics found
that every saved graph has at least $218$ low-defect five-sets whose
optimal-cut restriction family meets the optimal-cut family of the
remainder.  A separate maximum-cut-oriented diagnostic found many such
witnesses per optimum cut in the saved corpus.

However, abundance alone does not prove the desired theorem.  Even for the
balanced blow-up $B_3$, conditioning on preservation does not necessarily
lower the mean defect count below the target.  Thus a proof requires more
than a lower bound on the number of preserving five-sets.

\section{Optimal-face projection formulation}

For a five-set $Y$, define
\[
\operatorname{Res}_Y(G)=
\{f|_{G-Y}: f\in\Opt(G)\}.
\]
Then the H0 preservation property is exactly an intersection statement.

\begin{lemma}[Optimal-face equivalence]
For a fixed five-set $Y$, the following are equivalent:
\begin{enumerate}
\item there exists $f\in\Opt(G)$ with $f|_{G-Y}\in\Opt(G-Y)$;
\item $\operatorname{Res}_Y(G)\cap\Opt(G-Y)\neq\varnothing$;
\item there exists $f\in\Opt(G)$ with $\Delta_f(Y)=0$.
\end{enumerate}
\end{lemma}

The full saved-corpus diagnostics verify that every saved graph has many
low-defect five-sets with nonempty intersection.  A direct projection
counting identity can be obtained by double counting pairs $(f,Y)$, but
naive averaging is obstructed because distinct remainder optima can have
very different numbers of extensions back to global optima.  Thus the
remaining proof problem is weighted rather than purely cardinality based.

\section{The extension frontier}

The same phenomenon has a second exact formulation that makes the
remainder/extension tradeoff explicit.  Fix $Y$ and put $H=G-Y$.
For each attainable integer $j\ge0$, define
\[
F_Y(j)=
\min\{\text{extension cost over }Y:
\text{the coloring of }H\text{ has }\dd(H)+j\text{ monochromatic edges}\},
\]
where the extension cost counts monochromatic edges in $G[Y]$ and between
$Y$ and $H$, minimized over the $2^5$ assignments on $Y$.

\begin{lemma}[Extension-frontier identity]
For every five-set $Y$,
\[
\boxed{
\dd(G)-\dd(G-Y)=\min_{j\ge0}\bigl(j+F_Y(j)\bigr).
}
\]
Moreover, $Y$ is optimal-restriction preserving if and only if the minimum
is attained at $j=0$, equivalently
\[
F_Y(0)=\dd(G)-\dd(G-Y).
\]
\end{lemma}

Thus the exact H0 condition is
\[
j+F_Y(j)\ge F_Y(0)
\qquad\text{for every attainable }j\ge1.
\]

The saved H0 witnesses show that this anchored combined-cost inequality is
robust even though much stronger frontier-shape hypotheses are false.
Among all $19{,}270$ saved witness frontiers, only $7{,}745$ satisfy the
tested adjacent finite-layer slope condition, and only $2{,}405$ have
$F_Y$ itself minimized at zero.  Saved witnesses can have large downward
adjacent jumps.  In particular, monotonicity, one-step Lipschitz behavior,
and discrete convexity of $F_Y$ are not universal.

The empirical statement that survives is weaker and more precise:
\[
\boxed{
\text{for some useful }Y,\quad j+F_Y(j)\text{ is minimized at }j=0.
}
\]

\section{Why several natural proof shortcuts fail}

The computational program deliberately tested and rejected several
stronger statements.

\paragraph{Fixed-set local certificates.}
Singleton, pair, star, biclique, and a 32-flip local system provide exact
identities and useful conditional bounds, but purely local five-set data
does not control global remainder reoptimization in general.

\paragraph{Simple witness selection rules.}
Feature mining over all saved five-sets found no universal simple rule such
as minimizing degree sum, boundary size, or a fixed induced five-vertex
type.  A strong two-statistic rule covered $19{,}133/19{,}270$ saved graphs
but still failed on $137$.

\paragraph{Hitting-set and uncrossing arguments.}
Minimal dangerous remainder flips can be taken connected, but can be large.
The relevant signed cut-gain function is not generally submodular, so
standard laminarity/uncrossing arguments do not directly force a
five-vertex transversal.

\paragraph{Lexicographic gain profiles.}
Refining the objective by the multiplicities of gain levels does not force
descent: isomorphic remainders can preserve the entire gain profile while
replacing maximizing supports.

\paragraph{Naive abundance averaging.}
There may be many preserving deletions, yet the preserved-family conditional
mean defect can exceed the target.  The obstruction is an extension-weight
bias: remainder optima can have different numbers of extensions to global
optima.

These failures are useful diagnostics.  They indicate that the remaining
problem is genuinely a global selection statement about the relation
between the optimal-cut families of $G$ and $G-Y$.

\section{Current theorem target}

The full saved-corpus evidence suggests the following global stability
statement.

\begin{conjecture}[Optimal-restriction five-set principle]
Let $G$ be triangle-free on $5k$ vertices.  Under the structural hypotheses
forced by a genuine minimal counterexample to the Erd\H{o}s bound, there
exists a five-set $Y$ and a maximum cut $f$ of $G$ such that
\[
f|_{G-Y}\in\Opt(G-Y)
\]
and
\[
R_f(Y)\le 2k-1.
\]
Equivalently,
\[
\dd(G)-\dd(G-Y)\le2k-1
\]
with no remainder reoptimization loss.
\end{conjecture}

A proof would close the five-set induction route immediately.  The current
exact blocker is to show that not every low-defect deletion pair can fail
optimal projection.  In optimal-face language, one seeks positivity of the
low-defect intersection count after correctly accounting for extension
multiplicities.  In frontier language, one seeks a five-set with
$F_Y(0)\le2k-1$ for which no positive remainder-cost layer beats $j=0$.

\section{Reproducibility and scope}

The project contains exact bitmask MaxCut code, deterministic full-cut
landscape audits, saved-corpus witness tables, independent validation tests,
and SHA-256 checksums for release artifacts.  The v3 global-cut analyses
were performed on the saved canonical $n=15$ corpus of $19{,}270$ graphs.
The principal new artifacts are the global-cut landscape summaries, H0
witness tables, extension-frontier diagnostics, and optimal-face projection
notes.

The following distinctions are maintained throughout:
\begin{enumerate}
\item exact symbolic identities;
\item exact finite computations on explicitly defined saved corpora;
\item conjectural general statements motivated by those computations.
\end{enumerate}
No claim is made that the saved $19{,}270$-graph corpus is the complete
universe of triangle-free graphs on $15$ vertices unless a separate
exhaustive-generation proof is supplied.

\section{Discussion}

The computational search changed direction substantially over the course
of the project.  Early work focused on increasingly rich local recoloring
certificates.  Exact counterexamples to several stronger local statements
showed that the difficult part is not the five deleted vertices themselves,
but the global reoptimization of the remainder.

The full saved-corpus H0 audit then exposed a cleaner global pattern: after
choosing the right five-set, the troublesome reoptimization term vanishes.
That phenomenon survived all $19{,}270$ saved graphs and all tested hard
subclasses, while many stronger regularity properties of the extension
frontier failed.  The current evidence therefore points toward a theorem
about projection/intersection of optimal-cut families rather than a theorem
about local flip geometry.

\section*{Status and non-claims}

The Erd\H{o}s triangle-free bipartization conjecture remains open.  This
manuscript does not claim a proof of the conjecture or of the five-set
induction target in general.  The $19{,}270$-graph H0 statement is an exact
saved-corpus result, not an exhaustive $n=15$ classification claim.  The
purpose of the v3 note is to document the global cut-stability phenomenon,
its exact equivalent formulations, and the strongest proof barriers
identified so far.

\begin{thebibliography}{99}
\bibitem{BCL}
J. Balogh, F. C. Clemen, and B. Lidick\'y,
\emph{Max Cuts in Triangle-Free Graphs},
Extended Abstracts EuroComb 2021, Trends in Mathematics, Springer, 2021;
arXiv:2103.14179.

\bibitem{Ferudun}
A. Ferudun,
\emph{The Erd\H{o}s $n^2/25$ max-cut conjecture for small multiples of five,
via a per-root-MaxCut envelope and blow-up integrality},
arXiv:2606.28041, 2026.

\bibitem{CappelloSteffen}
C. Cappello and E. Steffen,
\emph{Frustration-critical signed graphs},
Discrete Applied Mathematics 322 (2022), 183--193;
arXiv:2112.02664.

\bibitem{Schrijver}
A. Schrijver,
\emph{A short proof of Guenin's characterization of weakly bipartite graphs},
Journal of Combinatorial Theory, Series B 85 (2002), 79--83.

\bibitem{Guenin}
J. F. Geelen and B. Guenin,
\emph{Packing Odd Circuits in Eulerian Graphs},
Journal of Combinatorial Theory, Series B 86 (2002), 280--295.
\end{thebibliography}

\end{document}
